% ============================================================================
% Preámbulo del informe (LuaLaTeX). Tipografía: Libertinus (texto y matemáticas)
% e Inter (títulos), la misma que usan las figuras generadas.
% ============================================================================
\usepackage[spanish,es-nodecimaldot,es-tabla,es-noshorthands]{babel}
\usepackage{amsmath}
\usepackage{fontspec}
\usepackage{unicode-math}
\setmainfont{Libertinus Serif}
\setsansfont{Inter}[Scale=MatchLowercase]
\setmonofont{Latin Modern Mono}[Scale=MatchLowercase,BoldFont={Latin Modern Mono},BoldFeatures={FakeBold=2}]
\setmathfont{Libertinus Math}
\setmathfont{Libertinus Math}[version=bold,FakeBold=1.6]
\usepackage[a4paper,margin=2.6cm,headheight=14pt]{geometry}
\usepackage{microtype}
\usepackage[dvipsnames]{xcolor}
\usepackage{graphicx}
\usepackage{booktabs}
\usepackage{tabularx}
\usepackage{xltabular}
\usepackage{array}
\usepackage{multirow}
\usepackage{caption}
\usepackage{float}
\usepackage[shortlabels]{enumitem}
\usepackage{csquotes}
\usepackage{fancyhdr}
\usepackage{titlesec}
\usepackage[most]{tcolorbox}
\usepackage{listings}
\usepackage{siunitx}
\usepackage[backend=biber,style=authoryear,maxcitenames=2,uniquelist=false,uniquename=false,dashed=false,useprefix=true]{biblatex}
\usepackage[hidelinks,pdfusetitle,bookmarksnumbered,hypertexnames=false]{hyperref}
\usepackage[spanish,noabbrev]{cleveref}

\addbibresource{bibliografia.bib}
\crefname{table}{tabla}{tablas}
\Crefname{table}{Tabla}{Tablas}
\crefname{lstlisting}{código}{códigos}
\Crefname{lstlisting}{Código}{Códigos}

% ------------------------------------------------------------------ colores
\definecolor{acento}{HTML}{2A78D6}
\definecolor{tinta2}{HTML}{52514E}
\definecolor{fondo}{HTML}{F3F2EE}

% ------------------------------------------------------------------ números
\sisetup{
  output-decimal-marker = {,},
  group-separator = {\,},
  group-minimum-digits = 5,
  detect-all,
  table-number-alignment = right,
}

% ---------------------------------------------------------------- resultados
% Cada cifra citada viene de la última corrida (generado/resultados.tex). Si una
% corrida posterior deja de producir una clave, el informe compila igualmente,
% la marca en rojo y emite una advertencia (el CI la convierte en error).
\makeatletter
\newcommand{\resdef}[2]{\expandafter\def\csname cs@res@#1\endcsname{#2}}
\newcommand{\res}[1]{%
  \ifcsname cs@res@#1\endcsname\csname cs@res@#1\endcsname
  \else\textcolor{red}{\textbf{[#1]}}\PackageWarning{resultados}{Falta la cifra «#1»}\fi}
\makeatother
% =====================================================================
% Generado automáticamente por cancerstats a partir de la corrida
% 20261009-095405-f8341612. No editar: se sobrescribe en cada exportación.
% =====================================================================
\resdef{aic-final}{\num{24653.2}}
\resdef{aj-medio-oeste}{\num{180.2}}
\resdef{aj-medio-oeste-ic}{$[\num{176.1};$\allowbreak\ $\num{184.2}]$}
\resdef{aj-noreste}{\num{168.8}}
\resdef{aj-noreste-ic}{$[\num{165.0};$\allowbreak\ $\num{172.7}]$}
\resdef{aj-oeste}{\num{170.4}}
\resdef{aj-oeste-ic}{$[\num{165.8};$\allowbreak\ $\num{175.0}]$}
\resdef{aj-sur}{\num{184.5}}
\resdef{aj-sur-ic}{$[\num{181.9};$\allowbreak\ $\num{187.2}]$}
\resdef{anova-estado-f}{\num{30.9}}
\resdef{anova-f}{\num{178.0}}
\resdef{anova-p}{$p < \num{0.001}$}
\resdef{asimetria}{\num{0.275}}
\resdef{b-bach1824}{\num{0.241}}
\resdef{b-bach25}{\num{0.354}}
\resdef{b-const}{\num{184.526}}
\resdef{b-desempleo}{\num{0.366}}
\resdef{b-desempleoxmedio-oeste}{\num{1.478}}
\resdef{b-desempleoxnoreste}{\num{0.903}}
\resdef{b-desempleoxoeste}{\num{-0.429}}
\resdef{b-edad}{\num{-0.788}}
\resdef{b-grado25}{\num{-1.054}}
\resdef{b-grado25xmedio-oeste}{\num{0.229}}
\resdef{b-grado25xnoreste}{\num{0.844}}
\resdef{b-grado25xoeste}{\num{-0.031}}
\resdef{b-hogar}{\num{-12.875}}
\resdef{b-incidencia}{\num{0.187}}
\resdef{b-medio-oeste}{\num{-4.327}}
\resdef{b-natalidad}{\num{-0.381}}
\resdef{b-noreste}{\num{-15.683}}
\resdef{b-oeste}{\num{-14.131}}
\resdef{b-otraraza}{\num{-0.660}}
\resdef{b-solopublico}{\num{0.892}}
\resdef{beta-bach1824}{\num{0.078}}
\resdef{beta-bach25}{\num{0.091}}
\resdef{beta-desempleo}{\num{0.045}}
\resdef{beta-edad}{\num{-0.147}}
\resdef{beta-grado25}{\num{-0.207}}
\resdef{beta-hogar}{\num{-0.115}}
\resdef{beta-incidencia}{\num{0.379}}
\resdef{beta-natalidad}{\num{-0.027}}
\resdef{beta-otraraza}{\num{-0.084}}
\resdef{beta-solopublico}{\num{0.197}}
\resdef{boot-bach1824}{$[\num{0.107};$\allowbreak\ $\num{0.364}]$}
\resdef{boot-bach25}{$[\num{-0.106};$\allowbreak\ $\num{0.688}]$}
\resdef{boot-desempleo}{$[\num{-0.209};$\allowbreak\ $\num{0.952}]$}
\resdef{boot-desempleoxmedio-oeste}{$[\num{0.463};$\allowbreak\ $\num{2.383}]$}
\resdef{boot-desempleoxnoreste}{$[\num{-0.436};$\allowbreak\ $\num{2.080}]$}
\resdef{boot-desempleoxoeste}{$[\num{-1.913};$\allowbreak\ $\num{1.024}]$}
\resdef{boot-edad}{$[\num{-1.156};$\allowbreak\ $\num{-0.452}]$}
\resdef{boot-grado25}{$[\num{-1.632};$\allowbreak\ $\num{-0.615}]$}
\resdef{boot-grado25xmedio-oeste}{$[\num{-0.398};$\allowbreak\ $\num{0.899}]$}
\resdef{boot-grado25xnoreste}{$[\num{0.048};$\allowbreak\ $\num{1.293}]$}
\resdef{boot-grado25xoeste}{$[\num{-0.495};$\allowbreak\ $\num{0.582}]$}
\resdef{boot-hogar}{$[\num{-18.688};$\allowbreak\ $\num{-5.950}]$}
\resdef{boot-incidencia}{$[\num{0.148};$\allowbreak\ $\num{0.225}]$}
\resdef{boot-medio-oeste}{$[\num{-9.286};$\allowbreak\ $\num{0.724}]$}
\resdef{boot-natalidad}{$[\num{-0.883};$\allowbreak\ $\num{0.087}]$}
\resdef{boot-noreste}{$[\num{-19.793};$\allowbreak\ $\num{-7.610}]$}
\resdef{boot-oeste}{$[\num{-20.095};$\allowbreak\ $\num{-7.808}]$}
\resdef{boot-otraraza}{$[\num{-1.245};$\allowbreak\ $\num{-0.348}]$}
\resdef{boot-reps}{\num{999}}
\resdef{boot-solopublico}{$[\num{0.499};$\allowbreak\ $\num{1.334}]$}
\resdef{boxcox-ic}{$[\num{0.536};$\allowbreak\ $\num{0.839}]$}
\resdef{boxcox-lambda}{\num{0.687}}
\resdef{boxcox-modelo}{\num{0.725}}
\resdef{boxcox-modelo-ic}{$[\num{0.600};$\allowbreak\ $\num{0.850}]$}
\resdef{bp-antes}{\num{262.7}}
\resdef{bp-antes-p}{$p < \num{0.001}$}
\resdef{bp-despues}{\num{73.0}}
\resdef{bp-despues-p}{$p < \num{0.001}$}
\resdef{bp-pob-antes}{\num{201.9}}
\resdef{bp-pob-antes-p}{$p < \num{0.001}$}
\resdef{bp-pob-despues}{\num{209.99}}
\resdef{bp-pob-despues-p}{$p < \num{0.001}$}
\resdef{brecha-ajustada}{\num{14.1}}
\resdef{brecha-bruta}{\num{31.8}}
\resdef{bruta-medio-oeste}{\num{176.7}}
\resdef{bruta-noreste}{\num{172.5}}
\resdef{bruta-oeste}{\num{157.1}}
\resdef{bruta-sur}{\num{189.0}}
\resdef{cambio-bach1824}{\num{35.2}}
\resdef{cambio-bach25}{\num{34.1}}
\resdef{cambio-desempleo}{\num{141.7}}
\resdef{cambio-edad}{\num{-57.3}}
\resdef{cambio-grado25}{\num{13.1}}
\resdef{cambio-hogar}{\num{-67.9}}
\resdef{cambio-medio-oeste}{\num{-14.0}}
\resdef{cambio-natalidad}{\num{-124.6}}
\resdef{cambio-noreste}{\num{40.7}}
\resdef{cambio-oeste}{\num{-22.0}}
\resdef{cambio-otraraza}{\num{-70.9}}
\resdef{cambio-solopublico}{\num{-0.9}}
\resdef{centinela-incidencia}{\num{453.5494221}}
\resdef{chi-estado-empleo-p}{$p = \num{0.591}$}
\resdef{chi-estado-incidencia-p}{$p < \num{0.001}$}
\resdef{chi-estado-soloprivado-p}{$p = \num{0.052}$}
\resdef{cmp-final-aic}{\num{23433.3}}
\resdef{cmp-final-bic}{\num{23551.3}}
\resdef{cmp-final-f}{\num{6.85}}
\resdef{cmp-final-f-rob}{\num{5.71}}
\resdef{cmp-final-gl}{\num{6}}
\resdef{cmp-final-k}{\num{19}}
\resdef{cmp-final-p}{$p < \num{0.001}$}
\resdef{cmp-final-p-rob}{$p < \num{0.001}$}
\resdef{cmp-final-r2aj}{\num{0.543}}
\resdef{cmp-maximo-aic}{\num{23436.5}}
\resdef{cmp-maximo-bic}{\num{23584.0}}
\resdef{cmp-maximo-f}{\num{110.70}}
\resdef{cmp-maximo-f-rob}{\num{106.52}}
\resdef{cmp-maximo-gl}{\num{21}}
\resdef{cmp-maximo-k}{\num{24}}
\resdef{cmp-maximo-p}{$p < \num{0.001}$}
\resdef{cmp-maximo-p-rob}{$p < \num{0.001}$}
\resdef{cmp-maximo-r2aj}{\num{0.544}}
\resdef{cmp-nulo-aic}{\num{25083.0}}
\resdef{cmp-nulo-bic}{\num{25106.6}}
\resdef{cmp-nulo-k}{\num{3}}
\resdef{cmp-nulo-r2aj}{\num{0.154}}
\resdef{cmp-seleccionado-aic}{\num{23462.4}}
\resdef{cmp-seleccionado-bic}{\num{23545.0}}
\resdef{cmp-seleccionado-f}{\num{4.35}}
\resdef{cmp-seleccionado-f-rob}{\num{2.13}}
\resdef{cmp-seleccionado-gl}{\num{11}}
\resdef{cmp-seleccionado-k}{\num{13}}
\resdef{cmp-seleccionado-p}{$p < \num{0.001}$}
\resdef{cmp-seleccionado-p-rob}{$p = \num{0.036}$}
\resdef{cmp-seleccionado-r2aj}{\num{0.537}}
\resdef{cobertura-mediana}{\num{95.4}}
\resdef{codificacion}{utf-8}
\resdef{comunes-estrategias}{bachillerato (18–24), grado universitario (25+), incidencia de cáncer, otra raza, región censal y solo seguro público}
\resdef{condados-mir}{Aleutians West Census Area, Alaska}
\resdef{conf-anchura-norm-gra}{\num{44.5}}
\resdef{conf-anchura-norm-med}{\num{51.1}}
\resdef{conf-anchura-norm-peq}{\num{71.4}}
\resdef{conf-anchura-std}{\num{60.0}}
\resdef{conf-nominal}{\num{90}}
\resdef{conf-norm}{\num{90.8}}
\resdef{conf-norm-gra}{\num{92.6}}
\resdef{conf-norm-med}{\num{91.6}}
\resdef{conf-norm-peq}{\num{88.2}}
\resdef{conf-std}{\num{91.1}}
\resdef{conf-std-gra}{\num{96.1}}
\resdef{conf-std-med}{\num{95.6}}
\resdef{conf-std-peq}{\num{81.9}}
\resdef{confusion-bach25}{\num{27.3}}
\resdef{confusion-bach25-en}{grado universitario (25+)}
\resdef{confusoras}{bachillerato (25+)}
\resdef{corrida-commit}{\texttt{aa76b3794edd}}
\resdef{corrida-duracion}{\num{160}}
\resdef{corrida-fecha}{2026-10-09 09:56:44 UTC}
\resdef{corrida-huella}{\texttt{f8341612ab08}}
\resdef{corrida-id}{\texttt{20261009-095405-f8341612}}
\resdef{corrida-semilla}{\num{20261007}}
\resdef{curtosis}{\num{1.350}}
\resdef{curtosis-antes}{\num{2.95}}
\resdef{curvatura}{ninguna}
\resdef{cv-r2-bosque}{\num{0.536}}
\resdef{cv-r2-lasso}{\num{0.552}}
\resdef{cv-r2-mco-completo}{\num{0.551}}
\resdef{cv-r2-mco-efectos}{\num{0.527}}
\resdef{cv-r2-potenciacion}{\num{0.532}}
\resdef{cv-r2-red-elastica}{\num{0.552}}
\resdef{cv-r2-referencia}{\num{-0.003}}
\resdef{cv-r2-ridge}{\num{0.553}}
\resdef{cv-rmse-bosque}{\num{18.95}}
\resdef{cv-rmse-bosque-dt}{\num{1.01}}
\resdef{cv-rmse-lasso}{\num{18.60}}
\resdef{cv-rmse-lasso-dt}{\num{1.36}}
\resdef{cv-rmse-mco-completo}{\num{18.61}}
\resdef{cv-rmse-mco-completo-dt}{\num{1.37}}
\resdef{cv-rmse-mco-efectos}{\num{19.11}}
\resdef{cv-rmse-mco-efectos-dt}{\num{1.44}}
\resdef{cv-rmse-potenciacion}{\num{19.03}}
\resdef{cv-rmse-potenciacion-dt}{\num{0.90}}
\resdef{cv-rmse-red-elastica}{\num{18.60}}
\resdef{cv-rmse-red-elastica-dt}{\num{1.36}}
\resdef{cv-rmse-referencia}{\num{27.90}}
\resdef{cv-rmse-referencia-dt}{\num{1.11}}
\resdef{cv-rmse-ridge}{\num{18.59}}
\resdef{cv-rmse-ridge-dt}{\num{1.37}}
\resdef{dt}{\num{27.75}}
\resdef{dt-ic}{$[\num{27.07};$\allowbreak\ $\num{28.47}]$}
\resdef{dt-razon-edad}{\num{0.067}}
\resdef{dw-antes}{\num{1.828}}
\resdef{dw-final}{\num{1.878}}
\resdef{ee-bach1824}{\num{0.064}}
\resdef{ee-bach25}{\num{0.192}}
\resdef{ee-const}{\num{1.319}}
\resdef{ee-desempleo}{\num{0.284}}
\resdef{ee-desempleoxmedio-oeste}{\num{0.446}}
\resdef{ee-desempleoxnoreste}{\num{0.555}}
\resdef{ee-desempleoxoeste}{\num{0.688}}
\resdef{ee-edad}{\num{0.187}}
\resdef{ee-grado25}{\num{0.255}}
\resdef{ee-grado25xmedio-oeste}{\num{0.308}}
\resdef{ee-grado25xnoreste}{\num{0.219}}
\resdef{ee-grado25xoeste}{\num{0.257}}
\resdef{ee-hogar}{\num{3.292}}
\resdef{ee-incidencia}{\num{0.019}}
\resdef{ee-medio-oeste}{\num{2.419}}
\resdef{ee-natalidad}{\num{0.249}}
\resdef{ee-noreste}{\num{2.316}}
\resdef{ee-oeste}{\num{2.756}}
\resdef{ee-otraraza}{\num{0.210}}
\resdef{ee-solopublico}{\num{0.214}}
\resdef{efiqr-bach1824}{\num{2.8}}
\resdef{efiqr-bach1824-ic}{$[\num{1.3};$\allowbreak\ $\num{4.3}]$}
\resdef{efiqr-bach25}{\num{3.3}}
\resdef{efiqr-bach25-ic}{$[\num{-0.3};$\allowbreak\ $\num{6.8}]$}
\resdef{efiqr-desempleo-medio-oeste}{\num{7.7}}
\resdef{efiqr-desempleo-noreste}{\num{5.3}}
\resdef{efiqr-desempleo-oeste}{\num{-0.3}}
\resdef{efiqr-desempleo-sur}{\num{1.5}}
\resdef{efiqr-edad}{\num{-4.8}}
\resdef{efiqr-edad-ic}{$[\num{-7.1};$\allowbreak\ $\num{-2.5}]$}
\resdef{efiqr-grado25-medio-oeste}{\num{-5.5}}
\resdef{efiqr-grado25-noreste}{\num{-1.4}}
\resdef{efiqr-grado25-oeste}{\num{-7.3}}
\resdef{efiqr-grado25-sur}{\num{-7.1}}
\resdef{efiqr-hogar}{\num{-3.3}}
\resdef{efiqr-hogar-ic}{$[\num{-5.1};$\allowbreak\ $\num{-1.6}]$}
\resdef{efiqr-incidencia}{\num{12.4}}
\resdef{efiqr-incidencia-ic}{$[\num{9.8};$\allowbreak\ $\num{15.0}]$}
\resdef{efiqr-natalidad}{\num{-0.8}}
\resdef{efiqr-natalidad-ic}{$[\num{-1.7};$\allowbreak\ $\num{0.2}]$}
\resdef{efiqr-otraraza}{\num{-1.2}}
\resdef{efiqr-otraraza-ic}{$[\num{-2.0};$\allowbreak\ $\num{-0.4}]$}
\resdef{efiqr-solopublico}{\num{7.4}}
\resdef{efiqr-solopublico-ic}{$[\num{3.8};$\allowbreak\ $\num{10.9}]$}
\resdef{eliminadas}{pobreza, log(población), población asiática, sin bachillerato (18–24), empleo (16+), log(1 + ensayos per cápita), población negra, hogares de matrimonios, grado universitario (18–24), bachillerato (25+), población nativa o multirracial y seguro privado vía empleador}
\resdef{ensayos-dif}{\num{-6.2}}
\resdef{ensayos-mw-p}{$p < \num{0.001}$}
\resdef{ensayos-n-con}{\num{1116}}
\resdef{ensayos-welch-p}{$p < \num{0.001}$}
\resdef{error-identidad}{\num{0.10}}
\resdef{estados-centinela}{Kansas, Minnesota y Nevada}
\resdef{eta2-estado}{\num{0.336}}
\resdef{eta2-region}{\num{0.149}}
\resdef{f-clasico}{\num{187.3}}
\resdef{f-robusto}{\num{94.3}}
\resdef{f-robusto-p}{$p < \num{0.001}$}
\resdef{fiv-max-antes}{\num{25.9}}
\resdef{fiv-max-despues}{\num{6.88}}
\resdef{fiv-max-final}{\num{6.25}}
\resdef{fiv-poda-casados}{\num{14.6}}
\resdef{fiv-poda-privado}{\num{13.9}}
\resdef{fiv-poda-publico}{\num{25.9}}
\resdef{fiv-poda-renta}{\num{11.0}}
\resdef{frac-muestreo-10000}{\num{52}}
\resdef{frac-muestreo-100000}{\num{10}}
\resdef{frac-muestreo-1000000}{\num{1}}
\resdef{frac-muestreo-2000}{\num{85}}
\resdef{frac-muestreo-25000}{\num{30}}
\resdef{gcv-r2-bosque}{\num{0.411}}
\resdef{gcv-r2-lasso}{\num{0.424}}
\resdef{gcv-r2-mco-completo}{\num{0.369}}
\resdef{gcv-r2-mco-efectos}{\num{0.435}}
\resdef{gcv-r2-potenciacion}{\num{0.414}}
\resdef{gcv-r2-red-elastica}{\num{0.424}}
\resdef{gcv-r2-referencia}{\num{-0.139}}
\resdef{gcv-r2-ridge}{\num{0.420}}
\resdef{gcv-rmse-bosque}{\num{19.93}}
\resdef{gcv-rmse-lasso}{\num{19.65}}
\resdef{gcv-rmse-mco-completo}{\num{20.36}}
\resdef{gcv-rmse-mco-efectos}{\num{19.43}}
\resdef{gcv-rmse-potenciacion}{\num{19.81}}
\resdef{gcv-rmse-red-elastica}{\num{19.65}}
\resdef{gcv-rmse-referencia}{\num{27.93}}
\resdef{gcv-rmse-ridge}{\num{19.68}}
\resdef{gl-inferencia}{\num{47}}
\resdef{gq-despues}{\num{3.71}}
\resdef{huber}{\num{178.26}}
\resdef{ic-bach1824}{$[\num{0.112};$\allowbreak\ $\num{0.370}]$}
\resdef{ic-bach25}{$[\num{-0.031};$\allowbreak\ $\num{0.739}]$}
\resdef{ic-const}{$[\num{181.873};$\allowbreak\ $\num{187.178}]$}
\resdef{ic-desempleo}{$[\num{-0.205};$\allowbreak\ $\num{0.937}]$}
\resdef{ic-desempleoxmedio-oeste}{$[\num{0.581};$\allowbreak\ $\num{2.375}]$}
\resdef{ic-desempleoxnoreste}{$[\num{-0.214};$\allowbreak\ $\num{2.020}]$}
\resdef{ic-desempleoxoeste}{$[\num{-1.813};$\allowbreak\ $\num{0.955}]$}
\resdef{ic-edad}{$[\num{-1.164};$\allowbreak\ $\num{-0.413}]$}
\resdef{ic-grado25}{$[\num{-1.566};$\allowbreak\ $\num{-0.542}]$}
\resdef{ic-grado25xmedio-oeste}{$[\num{-0.391};$\allowbreak\ $\num{0.849}]$}
\resdef{ic-grado25xnoreste}{$[\num{0.404};$\allowbreak\ $\num{1.285}]$}
\resdef{ic-grado25xoeste}{$[\num{-0.547};$\allowbreak\ $\num{0.485}]$}
\resdef{ic-hogar}{$[\num{-19.499};$\allowbreak\ $\num{-6.252}]$}
\resdef{ic-incidencia}{$[\num{0.148};$\allowbreak\ $\num{0.226}]$}
\resdef{ic-medio-oeste}{$[\num{-9.194};$\allowbreak\ $\num{0.541}]$}
\resdef{ic-natalidad}{$[\num{-0.883};$\allowbreak\ $\num{0.120}]$}
\resdef{ic-noreste}{$[\num{-20.343};$\allowbreak\ $\num{-11.023}]$}
\resdef{ic-oeste}{$[\num{-19.675};$\allowbreak\ $\num{-8.586}]$}
\resdef{ic-otraraza}{$[\num{-1.082};$\allowbreak\ $\num{-0.238}]$}
\resdef{ic-solopublico}{$[\num{0.463};$\allowbreak\ $\num{1.322}]$}
\resdef{icc-antes}{\num{0.074}}
\resdef{icc-antes-f}{\num{5.41}}
\resdef{icc-antes-p}{$p < \num{0.001}$}
\resdef{icc-final}{\num{0.079}}
\resdef{icc-final-f}{\num{6.01}}
\resdef{icc-final-p}{$p < \num{0.001}$}
\resdef{icc-mixto}{\num{0.118}}
\resdef{importancia1}{incidencia de cáncer}
\resdef{importancia2}{solo seguro público}
\resdef{importancia3}{hogares de matrimonios}
\resdef{influyente1}{Williamsburg city, Virginia}
\resdef{influyente1-cook}{\num{0.098}}
\resdef{influyente1-palanca}{\num{0.051}}
\resdef{ingenua-extra}{bachillerato (25+), seguro privado vía empleador, población negra, población nativa o multirracial y hogares de matrimonios}
\resdef{int-desempleoxregion-f}{\num{4.22}}
\resdef{int-desempleoxregion-p}{$p = \num{0.010}$}
\resdef{int-desempleoxregion-padj}{$p = \num{0.040}$}
\resdef{int-grado25xregion-f}{\num{7.27}}
\resdef{int-grado25xregion-p}{$p < \num{0.001}$}
\resdef{int-grado25xregion-padj}{$p = \num{0.002}$}
\resdef{interacciones}{Grado universitario (25+) $\times$ región y Desempleo (16+) $\times$ región}
\resdef{iqr-bach1824}{\num{11.50}}
\resdef{iqr-bach25}{\num{9.25}}
\resdef{iqr-desempleo}{\num{4.20}}
\resdef{iqr-edad}{\num{6.10}}
\resdef{iqr-grado25}{\num{6.70}}
\resdef{iqr-hogar}{\num{0.26}}
\resdef{iqr-incidencia}{\num{66.40}}
\resdef{iqr-natalidad}{\num{1.97}}
\resdef{iqr-otraraza}{\num{1.88}}
\resdef{iqr-solopublico}{\num{8.25}}
\resdef{k-final}{\num{19}}
\resdef{k-maximo}{\num{22}}
\resdef{kruskal-h}{\num{478.3}}
\resdef{kruskal-p}{$p < \num{0.001}$}
\resdef{lasso-nocero}{\num{65}}
\resdef{lasso-total}{\num{82}}
\resdef{levene-est}{\num{24.79}}
\resdef{levene-p}{$p < \num{0.001}$}
\resdef{little-aleatorio-est}{\num{62.8}}
\resdef{little-aleatorio-gl}{\num{74}}
\resdef{little-aleatorio-p}{$p = \num{0.820}$}
\resdef{little-aleatorio-patrones}{\num{4}}
\resdef{little-todo-est}{\num{522.3}}
\resdef{little-todo-gl}{\num{177}}
\resdef{little-todo-p}{$p < \num{0.001}$}
\resdef{little-todo-patrones}{\num{8}}
\resdef{m-imputaciones}{\num{20}}
\resdef{max-condado}{Union County, Florida}
\resdef{max-valor}{\num{362.8}}
\resdef{mayor-residuo}{Madison County, Mississippi}
\resdef{mayor-residuo-ajuste}{\num{158.4}}
\resdef{mayor-residuo-valor}{\num{7.35}}
\resdef{mayor-residuo-y}{\num{292.5}}
\resdef{mcd-corte}{\num{40.6}}
\resdef{mcd-p}{\num{25}}
\resdef{meda}{\num{17.0}}
\resdef{media}{\num{178.66}}
\resdef{media-ic}{$[\num{177.68};$\allowbreak\ $\num{179.65}]$}
\resdef{media-medio-oeste}{\num{174.5}}
\resdef{media-medio-oeste-ic}{$[\num{173.0};$\allowbreak\ $\num{176.0}]$}
\resdef{media-noreste}{\num{172.5}}
\resdef{media-noreste-ic}{$[\num{170.6};$\allowbreak\ $\num{174.4}]$}
\resdef{media-oeste}{\num{158.0}}
\resdef{media-oeste-ic}{$[\num{155.4};$\allowbreak\ $\num{160.7}]$}
\resdef{media-residuos}{\num{-1.6e-13}}
\resdef{media-sur}{\num{189.0}}
\resdef{media-sur-ic}{$[\num{187.5};$\allowbreak\ $\num{190.4}]$}
\resdef{mediana}{\num{178.1}}
\resdef{mediana-ic}{$[\num{177.0};$\allowbreak\ $\num{179.1}]$}
\resdef{mediana-ic-boot}{$[\num{177.0};$\allowbreak\ $\num{179.1}]$}
\resdef{min-condado}{Pitkin County, Colorado}
\resdef{min-valor}{\num{59.7}}
\resdef{mir-mediana}{\num{0.40}}
\resdef{modelo-1ee}{lasso}
\resdef{modelo-elegido}{ridge}
\resdef{modelo-mejor-cv}{ridge}
\resdef{modelo-mejor-gcv}{MCO con la especificación del modelo de efectos}
\resdef{n-atipicos-respuesta}{\num{18}}
\resdef{n-ausente-empleo}{\num{152}}
\resdef{n-ausente-incidencia}{\num{207}}
\resdef{n-ausente-soloprivado}{\num{609}}
\resdef{n-candidatas}{\num{25}}
\resdef{n-centinela}{\num{206}}
\resdef{n-columnas}{\num{35}}
\resdef{n-completos-todo}{\num{2175}}
\resdef{n-comunes-estrategias}{\num{6}}
\resdef{n-condados}{\num{3047}}
\resdef{n-confusoras}{\num{1}}
\resdef{n-cook}{\num{185}}
\resdef{n-cook-f50}{\num{0}}
\resdef{n-corr}{\num{25}}
\resdef{n-corr-sig}{\num{22}}
\resdef{n-curvatura}{\num{0}}
\resdef{n-edad-meses}{\num{30}}
\resdef{n-entrenamiento}{\num{2437}}
\resdef{n-especificaciones}{\num{10}}
\resdef{n-estados-final}{\num{48}}
\resdef{n-estudentizados}{\num{29}}
\resdef{n-final}{\num{2840}}
\resdef{n-hogar}{\num{61}}
\resdef{n-huber}{\num{618}}
\resdef{n-identidad-completos}{\num{762}}
\resdef{n-imputacion}{\num{3047}}
\resdef{n-ingenua}{\num{15}}
\resdef{n-ingenua-extra}{\num{5}}
\resdef{n-interacciones}{\num{2}}
\resdef{n-interacciones-candidatas}{\num{6}}
\resdef{n-mahal-clasica}{\num{330}}
\resdef{n-maximo}{\num{2699}}
\resdef{n-mcd}{\num{1202}}
\resdef{n-medio-oeste}{\num{1029}}
\resdef{n-mir}{\num{1}}
\resdef{n-nativa}{\num{86}}
\resdef{n-noreste}{\num{217}}
\resdef{n-oeste}{\num{418}}
\resdef{n-palanca}{\num{276}}
\resdef{n-pares-corr}{\num{6}}
\resdef{n-pasos-seleccion}{\num{12}}
\resdef{n-podadas}{\num{4}}
\resdef{n-prueba}{\num{610}}
\resdef{n-reglas}{\num{19}}
\resdef{n-reglas-advertencia}{\num{1}}
\resdef{n-reglas-error}{\num{4}}
\resdef{n-reglas-error-dep}{\num{0}}
\resdef{n-reglas-ok}{\num{12}}
\resdef{n-seleccionadas}{\num{10}}
\resdef{n-sin-influyentes}{\num{2655}}
\resdef{n-somecol}{\num{2285}}
\resdef{n-sur}{\num{1383}}
\resdef{n-terminos-final}{\num{13}}
\resdef{n-tukey-respuesta}{\num{64}}
\resdef{ncond-antes}{\num{972.2}}
\resdef{ncond-despues}{\num{187.1}}
\resdef{ncond-final}{\num{6.0}}
\resdef{norm-ad-est}{\num{4.074}}
\resdef{norm-ad-p}{$p = \num{0.010}$}
\resdef{norm-chi-est}{\num{94.270}}
\resdef{norm-chi-p}{$p < \num{0.001}$}
\resdef{norm-dagostino-est}{\num{127.363}}
\resdef{norm-dagostino-p}{$p < \num{0.001}$}
\resdef{norm-jb-est}{\num{269.820}}
\resdef{norm-jb-p}{$p < \num{0.001}$}
\resdef{norm-lilliefors-est}{\num{0.028}}
\resdef{norm-lilliefors-p}{$p = \num{0.001}$}
\resdef{norm-sw-est}{\num{0.990}}
\resdef{norm-sw-p}{$p < \num{0.001}$}
\resdef{p-bach1824}{$p < \num{0.001}$}
\resdef{p-bach25}{$p = \num{0.071}$}
\resdef{p-const}{$p < \num{0.001}$}
\resdef{p-desempleo}{$p = \num{0.204}$}
\resdef{p-desempleoxmedio-oeste}{$p = \num{0.002}$}
\resdef{p-desempleoxnoreste}{$p = \num{0.110}$}
\resdef{p-desempleoxoeste}{$p = \num{0.536}$}
\resdef{p-edad}{$p < \num{0.001}$}
\resdef{p-grado25}{$p < \num{0.001}$}
\resdef{p-grado25xmedio-oeste}{$p = \num{0.461}$}
\resdef{p-grado25xnoreste}{$p < \num{0.001}$}
\resdef{p-grado25xoeste}{$p = \num{0.905}$}
\resdef{p-hogar}{$p < \num{0.001}$}
\resdef{p-incidencia}{$p < \num{0.001}$}
\resdef{p-medio-oeste}{$p = \num{0.080}$}
\resdef{p-mw-mo-no}{$p = \num{0.263}$}
\resdef{p-natalidad}{$p = \num{0.132}$}
\resdef{p-noreste}{$p < \num{0.001}$}
\resdef{p-oeste}{$p < \num{0.001}$}
\resdef{p-otraraza}{$p = \num{0.003}$}
\resdef{p-solopublico}{$p < \num{0.001}$}
\resdef{p-tukey-mo-no}{$p = \num{0.731}$}
\resdef{pct-ausente-empleo}{\num{5.0}}
\resdef{pct-ausente-incidencia}{\num{6.8}}
\resdef{pct-ausente-soloprivado}{\num{20.0}}
\resdef{pct-brecha-explicada}{\num{56}}
\resdef{pct-completos-todo}{\num{71.4}}
\resdef{pct-estado}{\num{33.6}}
\resdef{pct-r2-final}{\num{55.8}}
\resdef{pct-region}{\num{14.9}}
\resdef{pct-somecol}{\num{75.0}}
\resdef{pend-desempleo-medio-oeste}{\num{1.844}}
\resdef{pend-desempleo-medio-oeste-ic}{$[\num{0.940};$\allowbreak\ $\num{2.749}]$}
\resdef{pend-desempleo-medio-oeste-p}{$p < \num{0.001}$}
\resdef{pend-desempleo-noreste}{\num{1.269}}
\resdef{pend-desempleo-noreste-ic}{$[\num{0.221};$\allowbreak\ $\num{2.317}]$}
\resdef{pend-desempleo-noreste-p}{$p = \num{0.019}$}
\resdef{pend-desempleo-oeste}{\num{-0.063}}
\resdef{pend-desempleo-oeste-ic}{$[\num{-1.417};$\allowbreak\ $\num{1.291}]$}
\resdef{pend-desempleo-oeste-p}{$p = \num{0.926}$}
\resdef{pend-desempleo-sur}{\num{0.366}}
\resdef{pend-desempleo-sur-ic}{$[\num{-0.205};$\allowbreak\ $\num{0.937}]$}
\resdef{pend-desempleo-sur-p}{$p = \num{0.204}$}
\resdef{pend-grado25-medio-oeste}{\num{-0.825}}
\resdef{pend-grado25-medio-oeste-ic}{$[\num{-1.542};$\allowbreak\ $\num{-0.108}]$}
\resdef{pend-grado25-medio-oeste-p}{$p = \num{0.025}$}
\resdef{pend-grado25-noreste}{\num{-0.209}}
\resdef{pend-grado25-noreste-ic}{$[\num{-0.723};$\allowbreak\ $\num{0.304}]$}
\resdef{pend-grado25-noreste-p}{$p = \num{0.416}$}
\resdef{pend-grado25-oeste}{\num{-1.085}}
\resdef{pend-grado25-oeste-ic}{$[\num{-1.600};$\allowbreak\ $\num{-0.569}]$}
\resdef{pend-grado25-oeste-p}{$p < \num{0.001}$}
\resdef{pend-grado25-sur}{\num{-1.054}}
\resdef{pend-grado25-sur-ic}{$[\num{-1.566};$\allowbreak\ $\num{-0.542}]$}
\resdef{pend-grado25-sur-p}{$p < \num{0.001}$}
\resdef{pobreza-d}{\num{0.78}}
\resdef{pobreza-dif}{\num{20.4}}
\resdef{pobreza-dif-ic}{$[\num{18.1};$\allowbreak\ $\num{22.8}]$}
\resdef{pobreza-f-p}{$p < \num{0.001}$}
\resdef{pobreza-ks}{\num{0.348}}
\resdef{pobreza-media-alta}{\num{193.5}}
\resdef{pobreza-media-baja}{\num{173.1}}
\resdef{pobreza-n-alta}{\num{830}}
\resdef{pobreza-welch-t}{\num{17.24}}
\resdef{podadas}{seguro público, población casada, seguro privado y log(renta mediana)}
\resdef{r-asiatica}{\num{-0.186}}
\resdef{r-asiatica-ic}{$[\num{-0.220};$\allowbreak\ $\num{-0.152}]$}
\resdef{r-bach1824}{\num{0.262}}
\resdef{r-bach1824-ic}{$[\num{0.229};$\allowbreak\ $\num{0.295}]$}
\resdef{r-bach25}{\num{0.405}}
\resdef{r-bach25-ic}{$[\num{0.374};$\allowbreak\ $\num{0.434}]$}
\resdef{r-casados}{\num{-0.267}}
\resdef{r-casados-ic}{$[\num{-0.299};$\allowbreak\ $\num{-0.234}]$}
\resdef{r-desempleo}{\num{0.378}}
\resdef{r-desempleo-ic}{$[\num{0.348};$\allowbreak\ $\num{0.408}]$}
\resdef{r-edad}{\num{-0.005}}
\resdef{r-edad-ic}{$[\num{-0.040};$\allowbreak\ $\num{0.031}]$}
\resdef{r-empleo}{\num{-0.412}}
\resdef{r-empleo-ic}{$[\num{-0.442};$\allowbreak\ $\num{-0.381}]$}
\resdef{r-ensayos}{\num{-0.100}}
\resdef{r-ensayos-ic}{$[\num{-0.135};$\allowbreak\ $\num{-0.064}]$}
\resdef{r-grado1824}{\num{-0.288}}
\resdef{r-grado1824-ic}{$[\num{-0.320};$\allowbreak\ $\num{-0.255}]$}
\resdef{r-grado25}{\num{-0.485}}
\resdef{r-grado25-ic}{$[\num{-0.512};$\allowbreak\ $\num{-0.458}]$}
\resdef{r-hogar}{\num{-0.037}}
\resdef{r-hogar-ic}{$[\num{-0.073};$\allowbreak\ $\num{-0.002}]$}
\resdef{r-hogarescasados}{\num{-0.293}}
\resdef{r-hogarescasados-ic}{$[\num{-0.325};$\allowbreak\ $\num{-0.261}]$}
\resdef{r-incidencia}{\num{0.472}}
\resdef{r-incidencia-ic}{$[\num{0.443};$\allowbreak\ $\num{0.500}]$}
\resdef{r-natalidad}{\num{-0.087}}
\resdef{r-natalidad-ic}{$[\num{-0.123};$\allowbreak\ $\num{-0.052}]$}
\resdef{r-nativa}{\num{0.043}}
\resdef{r-nativa-ic}{$[\num{0.007};$\allowbreak\ $\num{0.078}]$}
\resdef{r-negra}{\num{0.257}}
\resdef{r-negra-ic}{$[\num{0.224};$\allowbreak\ $\num{0.290}]$}
\resdef{r-otraraza}{\num{-0.190}}
\resdef{r-otraraza-ic}{$[\num{-0.224};$\allowbreak\ $\num{-0.155}]$}
\resdef{r-poblacion}{\num{-0.071}}
\resdef{r-poblacion-ic}{$[\num{-0.106};$\allowbreak\ $\num{-0.035}]$}
\resdef{r-pobreza}{\num{0.429}}
\resdef{r-pobreza-ic}{$[\num{0.400};$\allowbreak\ $\num{0.458}]$}
\resdef{r-privado}{\num{-0.386}}
\resdef{r-privado-ic}{$[\num{-0.416};$\allowbreak\ $\num{-0.355}]$}
\resdef{r-privempleador}{\num{-0.267}}
\resdef{r-privempleador-ic}{$[\num{-0.300};$\allowbreak\ $\num{-0.234}]$}
\resdef{r-publico}{\num{0.405}}
\resdef{r-publico-ic}{$[\num{0.374};$\allowbreak\ $\num{0.434}]$}
\resdef{r-renta}{\num{-0.452}}
\resdef{r-renta-ic}{$[\num{-0.480};$\allowbreak\ $\num{-0.424}]$}
\resdef{r-sinbach1824}{\num{0.088}}
\resdef{r-sinbach1824-ic}{$[\num{0.053};$\allowbreak\ $\num{0.124}]$}
\resdef{r-solopublico}{\num{0.449}}
\resdef{r-solopublico-ic}{$[\num{0.421};$\allowbreak\ $\num{0.477}]$}
\resdef{r2-con-incidencia}{\num{0.558}}
\resdef{r2-final}{\num{0.558}}
\resdef{r2-inicial}{\num{0.548}}
\resdef{r2-noponderado}{\num{0.558}}
\resdef{r2-sin-incidencia}{\num{0.438}}
\resdef{r2aj-final}{\num{0.555}}
\resdef{rachas}{\num{1172}}
\resdef{rachas-esperadas}{\num{1522.5}}
\resdef{rachas-p}{$p < \num{0.001}$}
\resdef{rachas-res-p}{$p < \num{0.001}$}
\resdef{rachas-res-z}{\num{-3.33}}
\resdef{rachas-z}{\num{-12.71}}
\resdef{rango-bach1824}{$[\num{0.154};$\allowbreak\ $\num{0.271}]$}
\resdef{rango-bach25}{$[\num{0.151};$\allowbreak\ $\num{0.416}]$}
\resdef{rango-desempleo}{$[\num{0.283};$\allowbreak\ $\num{0.767}]$}
\resdef{rango-desempleoxmedio-oeste}{$[\num{0.760};$\allowbreak\ $\num{1.930}]$}
\resdef{rango-desempleoxnoreste}{$[\num{0.479};$\allowbreak\ $\num{0.988}]$}
\resdef{rango-desempleoxoeste}{$[\num{-1.380};$\allowbreak\ $\num{0.186}]$}
\resdef{rango-edad}{$[\num{-0.929};$\allowbreak\ $\num{-0.581}]$}
\resdef{rango-grado25}{$[\num{-1.283};$\allowbreak\ $\num{-1.015}]$}
\resdef{rango-grado25xmedio-oeste}{$[\num{0.032};$\allowbreak\ $\num{0.403}]$}
\resdef{rango-grado25xnoreste}{$[\num{0.599};$\allowbreak\ $\num{1.160}]$}
\resdef{rango-grado25xoeste}{$[\num{-0.094};$\allowbreak\ $\num{0.350}]$}
\resdef{rango-hogar}{$[\num{-16.176};$\allowbreak\ $\num{-10.647}]$}
\resdef{rango-incidencia}{$[\num{0.175};$\allowbreak\ $\num{0.188}]$}
\resdef{rango-medio-oeste}{$[\num{-5.910};$\allowbreak\ $\num{-4.302}]$}
\resdef{rango-natalidad}{$[\num{-0.587};$\allowbreak\ $\num{-0.356}]$}
\resdef{rango-noreste}{$[\num{-17.449};$\allowbreak\ $\num{-13.945}]$}
\resdef{rango-oeste}{$[\num{-17.063};$\allowbreak\ $\num{-13.341}]$}
\resdef{rango-otraraza}{$[\num{-0.769};$\allowbreak\ $\num{-0.548}]$}
\resdef{rango-solopublico}{$[\num{0.777};$\allowbreak\ $\num{0.935}]$}
\resdef{razon-edad}{\num{12.008}}
\resdef{razon-var-deciles}{\num{6.8}}
\resdef{recortada}{\num{178.26}}
\resdef{renta-decil1}{\num{202.7}}
\resdef{renta-decil10}{\num{159.1}}
\resdef{renta-spearman}{\num{-0.461}}
\resdef{res-asimetria}{\num{0.184}}
\resdef{res-curtosis}{\num{2.81}}
\resdef{res-jb}{\num{952.9}}
\resdef{res-jb-p}{$p < \num{0.001}$}
\resdef{res-sw}{\num{0.9750}}
\resdef{res-sw-p}{$p < \num{0.001}$}
\resdef{reset-f}{\num{1.291}}
\resdef{reset-p}{$p = \num{0.275}$}
\resdef{rho-asiatica}{\num{-0.206}}
\resdef{rho-bach1824}{\num{0.276}}
\resdef{rho-bach25}{\num{0.419}}
\resdef{rho-casados}{\num{-0.259}}
\resdef{rho-desempleo}{\num{0.403}}
\resdef{rho-edad}{\num{0.007}}
\resdef{rho-empleo}{\num{-0.432}}
\resdef{rho-ensayos}{\num{-0.113}}
\resdef{rho-grado1824}{\num{-0.271}}
\resdef{rho-grado25}{\num{-0.502}}
\resdef{rho-hogar}{\num{0.023}}
\resdef{rho-hogarescasados}{\num{-0.280}}
\resdef{rho-incidencia}{\num{0.428}}
\resdef{rho-natalidad}{\num{-0.053}}
\resdef{rho-nativa}{\num{-0.139}}
\resdef{rho-negra}{\num{0.255}}
\resdef{rho-otraraza}{\num{-0.215}}
\resdef{rho-poblacion}{\num{-0.044}}
\resdef{rho-pobreza}{\num{0.443}}
\resdef{rho-privado}{\num{-0.411}}
\resdef{rho-privempleador}{\num{-0.277}}
\resdef{rho-publico}{\num{0.401}}
\resdef{rho-renta}{\num{-0.463}}
\resdef{rho-sinbach1824}{\num{0.126}}
\resdef{rho-solopublico}{\num{0.453}}
\resdef{rmse-final}{\num{18.44}}
\resdef{seleccion-metodo}{hacia atrás}
\resdef{sha-corto}{\texttt{1924e0da423a7408}}
\resdef{sigma-final}{\num{18.50}}
\resdef{sin-inc-bach1824}{\num{0.326}}
\resdef{sin-inc-bach25}{\num{0.474}}
\resdef{sin-inc-desempleo}{\num{0.885}}
\resdef{sin-inc-edad}{\num{-1.240}}
\resdef{sin-inc-grado25}{\num{-0.916}}
\resdef{sin-inc-hogar}{\num{-21.614}}
\resdef{sin-inc-medio-oeste}{\num{-4.934}}
\resdef{sin-inc-natalidad}{\num{-0.857}}
\resdef{sin-inc-noreste}{\num{-9.303}}
\resdef{sin-inc-oeste}{\num{-17.237}}
\resdef{sin-inc-otraraza}{\num{-1.127}}
\resdef{sin-inc-solopublico}{\num{0.884}}
\resdef{t-bach1824}{\num{3.76}}
\resdef{t-bach25}{\num{1.85}}
\resdef{t-const}{\num{139.95}}
\resdef{t-desempleo}{\num{1.29}}
\resdef{t-desempleoxmedio-oeste}{\num{3.31}}
\resdef{t-desempleoxnoreste}{\num{1.63}}
\resdef{t-desempleoxoeste}{\num{-0.62}}
\resdef{t-edad}{\num{-4.22}}
\resdef{t-grado25}{\num{-4.14}}
\resdef{t-grado25xmedio-oeste}{\num{0.74}}
\resdef{t-grado25xnoreste}{\num{3.86}}
\resdef{t-grado25xoeste}{\num{-0.12}}
\resdef{t-hogar}{\num{-3.91}}
\resdef{t-incidencia}{\num{9.61}}
\resdef{t-medio-oeste}{\num{-1.79}}
\resdef{t-natalidad}{\num{-1.53}}
\resdef{t-noreste}{\num{-6.77}}
\resdef{t-oeste}{\num{-5.13}}
\resdef{t-otraraza}{\num{-3.14}}
\resdef{t-solopublico}{\num{4.18}}
\resdef{test-dif-frente}{mco-efectos}
\resdef{test-dif-ic}{$[\num{-1.26};$\allowbreak\ $\num{-0.40}]$}
\resdef{test-mae-bosque}{\num{13.22}}
\resdef{test-mae-elegido}{\num{13.18}}
\resdef{test-mae-lasso}{\num{13.17}}
\resdef{test-mae-mco-completo}{\num{13.26}}
\resdef{test-mae-mco-efectos}{\num{13.85}}
\resdef{test-mae-potenciacion}{\num{13.36}}
\resdef{test-mae-red-elastica}{\num{13.19}}
\resdef{test-mae-referencia}{\num{21.16}}
\resdef{test-mae-ridge}{\num{13.18}}
\resdef{test-r2-bosque}{\num{0.577}}
\resdef{test-r2-elegido}{\num{0.589}}
\resdef{test-r2-lasso}{\num{0.589}}
\resdef{test-r2-mco-completo}{\num{0.585}}
\resdef{test-r2-mco-efectos}{\num{0.549}}
\resdef{test-r2-potenciacion}{\num{0.577}}
\resdef{test-r2-red-elastica}{\num{0.589}}
\resdef{test-r2-referencia}{\num{-0.002}}
\resdef{test-r2-ridge}{\num{0.589}}
\resdef{test-rmse-bosque}{\num{17.60}}
\resdef{test-rmse-elegido}{\num{17.36}}
\resdef{test-rmse-ic}{$[\num{16.25};$\allowbreak\ $\num{18.51}]$}
\resdef{test-rmse-lasso}{\num{17.35}}
\resdef{test-rmse-mco-completo}{\num{17.44}}
\resdef{test-rmse-mco-efectos}{\num{18.18}}
\resdef{test-rmse-potenciacion}{\num{17.61}}
\resdef{test-rmse-red-elastica}{\num{17.36}}
\resdef{test-rmse-referencia}{\num{27.10}}
\resdef{test-rmse-ridge}{\num{17.36}}
\resdef{umbral-cook}{\num{0.00141}}
\resdef{umbral-dfbetas}{\num{0.0375}}
\resdef{umbral-palanca}{\num{0.0141}}
\resdef{var-decil1}{\num{758}}
\resdef{var-decil10}{\num{112}}
\resdef{varianza}{\num{770.1}}
\resdef{varianza-ic}{$[\num{732.9};$\allowbreak\ $\num{810.3}]$}
\resdef{version-matplotlib}{3.11.2}
\resdef{version-numpy}{2.5.3}
\resdef{version-pandas}{3.0.6}
\resdef{version-python}{3.12.3}
\resdef{version-scikitlearn}{1.9.1}
\resdef{version-scipy}{1.18.1}
\resdef{version-statsmodels}{0.15.0}
\resdef{vf-a}{\num{195.2}}
\resdef{vf-b}{\num{2.136}}
\resdef{vf-b-p}{$p < \num{0.001}$}
\resdef{vf-iter}{\num{5}}
\resdef{viol-r01}{\num{0}}
\resdef{viol-r01-dep}{\num{0}}
\resdef{viol-r02}{\num{0}}
\resdef{viol-r02-dep}{\num{0}}
\resdef{viol-r03}{\num{0}}
\resdef{viol-r03-dep}{\num{0}}
\resdef{viol-r04}{\num{206}}
\resdef{viol-r04-dep}{\num{0}}
\resdef{viol-r05}{\num{30}}
\resdef{viol-r05-dep}{\num{0}}
\resdef{viol-r06}{\num{30}}
\resdef{viol-r06-dep}{\num{0}}
\resdef{viol-r07}{\num{61}}
\resdef{viol-r07-dep}{\num{0}}
\resdef{viol-r08}{\num{0}}
\resdef{viol-r08-dep}{\num{0}}
\resdef{viol-r09}{\num{0}}
\resdef{viol-r09-dep}{\num{0}}
\resdef{viol-r10}{\num{0}}
\resdef{viol-r10-dep}{\num{0}}
\resdef{viol-r11}{\num{0}}
\resdef{viol-r11-dep}{\num{0}}
\resdef{viol-r12}{\num{0}}
\resdef{viol-r12-dep}{\num{0}}
\resdef{viol-r13}{\num{0}}
\resdef{viol-r13-dep}{\num{0}}
\resdef{viol-r14}{\num{0}}
\resdef{viol-r14-dep}{\num{0}}
\resdef{viol-r15}{\num{1}}
\resdef{viol-r15-dep}{\num{0}}
\resdef{viol-r16}{\num{0}}
\resdef{viol-r16-dep}{\num{0}}
\resdef{viol-r17}{\num{2504}}
\resdef{viol-r17-dep}{\num{872}}
\resdef{viol-r18}{\num{1931}}
\resdef{viol-r18-dep}{\num{1931}}
\resdef{viol-r19}{\num{0}}
\resdef{viol-r19-dep}{\num{0}}
\resdef{welch-f}{\num{164.3}}
\resdef{welch-gl2}{\num{898.0}}
\resdef{welch-p}{$p < \num{0.001}$}
\resdef{white-despues}{\num{34.99}}
\resdef{white-despues-p}{$p < \num{0.001}$}


% ------------------------------------------------------------------ títulos
\titleformat{\chapter}[display]
  {\sffamily\bfseries\Huge}
  {\color{acento}\large\MakeUppercase{\chaptertitlename}\ \thechapter}{6pt}{}
\titlespacing*{\chapter}{0pt}{0pt}{24pt}
\titleformat{\section}{\sffamily\bfseries\Large}{\color{acento}\thesection}{0.8em}{}
\titleformat{\subsection}{\sffamily\bfseries\large}{\color{acento}\thesubsection}{0.8em}{}
\titleformat{\subsubsection}{\sffamily\bfseries}{}{0pt}{}

\captionsetup{font=small,labelfont={bf,sf},format=hang,skip=6pt}
\captionsetup[table]{position=top}
\setlength{\parskip}{0.35em}
\setlength{\emergencystretch}{1.5em}
\setlist{itemsep=0.2em,topsep=0.3em}
\renewcommand{\arraystretch}{1.12}

% ------------------------------------------------------------------ cabeceras
\pagestyle{fancy}
\fancyhf{}
\fancyhead[L]{\sffamily\footnotesize\color{tinta2}\leftmark}
\fancyhead[R]{\sffamily\footnotesize\color{tinta2}Mortalidad por cáncer en EE.\,UU.}
\fancyfoot[C]{\sffamily\footnotesize\thepage}
\renewcommand{\headrulewidth}{0.3pt}
\fancypagestyle{plain}{\fancyhf{}\fancyfoot[C]{\sffamily\footnotesize\thepage}\renewcommand{\headrulewidth}{0pt}}

% ------------------------------------------------------------------ recuadros
\newtcolorbox{idea}[1][]{enhanced,breakable,colback=acento!6,colframe=acento,
  boxrule=0pt,leftrule=2.5pt,arc=2pt,left=8pt,right=8pt,top=6pt,bottom=6pt,
  fonttitle=\sffamily\bfseries,coltitle=acento,attach title to upper={\ },title={#1}}
\newtcolorbox{entregable}[1]{enhanced,breakable,colback=fondo,colframe=fondo,
  boxrule=0pt,arc=3pt,left=8pt,right=8pt,top=6pt,bottom=6pt,
  fontupper=\small,title={\sffamily\bfseries Entregable #1},coltitle=black,
  colbacktitle=fondo,toptitle=2pt,bottomtitle=0pt}

% ------------------------------------------------------------------ código
\lstdefinestyle{codigo}{basicstyle=\ttfamily\footnotesize,columns=fullflexible,
  keepspaces=true,breaklines=true,frame=single,rulecolor=\color{black!15},
  backgroundcolor=\color{fondo},xleftmargin=4pt,xrightmargin=4pt,
  showstringspaces=false,upquote=true,
  literate={á}{{\'a}}1 {é}{{\'e}}1 {í}{{\'i}}1 {ó}{{\'o}}1 {ú}{{\'u}}1 {ñ}{{\~n}}1
           {Á}{{\'A}}1 {É}{{\'E}}1 {Í}{{\'I}}1 {Ó}{{\'O}}1 {Ú}{{\'U}}1 {Ñ}{{\~N}}1
           {«}{{\guillemotleft}}1 {»}{{\guillemotright}}1 {·}{{\textperiodcentered}}1
           {−}{{-}}1 {σ}{{$\sigma$}}1 {²}{{$^2$}}1 {≥}{{$\geq$}}1 {≤}{{$\leq$}}1}
\lstset{style=codigo}
\renewcommand{\lstlistingname}{Código}

% ------------------------------------------------------------------ utilidades
\newcommand{\var}[1]{\texttt{#1}}
\newcommand{\figura}[3][0.98]{%
  \begin{figure}[htbp]\centering
  \includegraphics[width=#1\linewidth]{generado/figuras/#2.pdf}
  \caption{#3}\label{fig:#2}\end{figure}}
\DeclareMathOperator{\Var}{Var}
\DeclareMathOperator{\Cov}{Cov}
\DeclareMathOperator{\E}{E}
\newcommand{\R}{\mathbb{R}}

% Las páginas que sólo contienen flotantes los colocan arriba, no centrados en vertical.
\makeatletter
\setlength{\@fptop}{0pt}
\setlength{\@fpsep}{12pt plus 1fil}
\setlength{\@fpbot}{0pt plus 1fil}
\makeatother
